\section{Large Language Models}

Large Language Models (LLMs) represent a paradigm shift in natural language processing, characterised by their ability to model, generate, and manipulate human language with remarkable fluency. At their core, LLMs are neural sequence models trained on vast corpora of textual data, typically employing the Transformer architecture introduced by Vaswani et al.\ \cite{VasShaPar:aiayn17}. The defining feature of this architecture is the self-attention mechanism, which enables the model to weigh the relevance of every token in the input sequence when computing a representation for any given position. Formally, for a sequence of input vectors \(\mathbf{X} = [\mathbf{x}_1, \mathbf{x}_2, \dots, \mathbf{x}_n]\), the scaled dot-product attention is computed as:

\[
\mathrm{Attention}(\mathbf{Q}, \mathbf{K}, \mathbf{V}) = \mathrm{softmax}\left(\frac{\mathbf{Q}\mathbf{K}^{\mathsf{T}}}{\sqrt{d_k}}\right)\mathbf{V},
\]

where \(\mathbf{Q}\), \(\mathbf{K}\), and \(\mathbf{V}\) are query, key, and value matrices derived from \(\mathbf{X}\), and \(d_k\) is the dimensionality of the key vectors. This mechanism allows each token to attend to all other tokens in the sequence, thereby capturing long-range dependencies that recurrent architectures struggle to model effectively.

\subsection{Architectural Evolution}

The original Transformer was proposed in an encoder--decoder configuration for machine translation. Subsequent work demonstrated that scaling the decoder-only variant---a causal language model that predicts the next token given all preceding tokens---yields emergent capabilities not explicitly programmed during training \cite{brown2020languagemodelsfewshotlearners}. The autoregressive objective can be expressed as:

\[
P(\mathbf{x}) = \prod_{t=1}^{T} P(x_t \mid x_{<t}; \theta),
\]

where \(\theta\) denotes the model parameters and \(x_{<t}\) represents the prefix of tokens preceding position \(t\). Models such as the Generative Pre-trained Transformer (GPT) family \cite{radford2018improving, brown2020languagemodelsfewshotlearners} and LLaMA \cite{touvron2023llama} adopt this formulation, scaling to hundreds of billions of parameters and training on trillions of tokens \cite{hoffmann2022training}.

\subsection{Pre-training and Scaling Laws}

LLMs are pre-trained in a self-supervised manner on large, diverse corpora drawn from the Web, books, and scientific articles. During pre-training, the model learns a rich distributional representation of language, capturing syntactic, semantic, and factual knowledge implicitly from the training data. Kaplan et al.\ \cite{kaplan2020scaling} and Hoffmann et al.\ \cite{hoffmann2022training} established that the performance of these models follows a power-law relationship with three factors: the number of model parameters, the size of the training dataset, and the amount of compute used for training. This finding, known as the \emph{scaling law}, has driven the trend toward increasingly large models and datasets.

\subsection{Instruction Tuning and Alignment}

While pre-trained LLMs exhibit strong language understanding, their outputs are not naturally aligned with user intent. To bridge this gap, \emph{instruction tuning} is employed: the model is fine-tuned on a dataset of (instruction, response) pairs, teaching it to follow diverse prompts \cite{ouyang2022training}. This process is often followed by Reinforcement Learning from Human Feedback (RLHF) \cite{ziegler2019fine, lambert2024illustrating}, wherein a reward model trained on human preferences is used to steer the policy toward outputs that are helpful, honest, and harmless. The combined effect of instruction tuning and RLHF is what makes models such as GPT-4 and Claude behave as conversational assistants rather than raw next-token predictors.

\subsection{Limitations}

Despite their capabilities, LLMs suffer from several well-documented limitations. First, they exhibit a tendency to \emph{hallucinate}---generating content that is fluent and superficially plausible but factually incorrect \cite{ji2023survey, huang2023survey}. Second, their knowledge is frozen at the time of pre-training, rendering them unable to incorporate new or domain-specific information without costly re-training or fine-tuning. Third, the parametric nature of their knowledge makes it difficult to verify the provenance of generated claims, as the model does not cite external sources by default. These limitations motivate the integration of LLMs with external retrieval mechanisms, a topic explored in the Retrieval-Augmented Generation section of this chapter.

\subsection{Chain-of-Thought Reasoning}

A notable advancement in eliciting sophisticated reasoning from LLMs is the \emph{chain-of-thought} (CoT) prompting technique, introduced by Wei et al.\ \cite{wei2022chain}. This method instructs the model to generate intermediate reasoning steps before arriving at a final answer, thereby decomposing complex multi-step problems into a sequence of manageable sub-tasks. Rather than producing a direct answer, the model is prompted to articulate a logical progression---for instance, by including the phrase ``Let us think step by step'' \cite{kojima2022large}---which mirrors the human cognitive process of deliberative reasoning.

Formally, given an input query \(q\) and a set of exemplars \(\mathcal{E} = \{(q_i, r_i, a_i)\}_{i=1}^{k}\), where each exemplar consists of a question \(q_i\), a reasoning chain \(r_i\), and an answer \(a_i\), the model generates a reasoning chain \(\hat{r}\) and a final answer \(\hat{a}\) autoregressively:

\[
\hat{r}, \hat{a} = \arg\max_{r,a} P(r, a \mid q, \mathcal{E}; \theta).
\]

The reasoning chain \(\hat{r}\) serves as an intermediate scaffold that guides the model toward the correct answer by reducing the cognitive load of a single-step prediction.

The effectiveness of CoT reasoning has been demonstrated across a wide range of arithmetic, commonsense, and symbolic reasoning tasks \cite{wei2022chain, wang2023selfconsistency}. Moreover, subsequent work has extended this paradigm in several directions. \emph{Self-consistency} \cite{wang2023selfconsistency} samples multiple reasoning paths and selects the most consistent answer via marginalisation, thereby improving robustness against stochastic variations in generated chains. \emph{Tree-of-Thoughts} (ToT) \cite{yao2023tree} generalises the linear chain into a tree structure, enabling the model to explore multiple reasoning branches and perform deliberate backtracking when a branch proves unfruitful. Similarly, \emph{Graph-of-Thoughts} (GoT) \cite{besta2024graph} further extends this concept by modelling reasoning steps as a directed graph, allowing for aggregation and refinement of partial solutions.

Despite its strengths, CoT prompting is not without limitations. The quality of the generated reasoning chain is contingent upon the model's underlying capabilities; if the model lacks the requisite knowledge, the chain may propagate errors rather than mitigate them. Furthermore, CoT incurs additional computational cost due to the generation of extended token sequences, and its benefits are most pronounced in tasks that genuinely require multi-step reasoning, while being marginal for simpler queries \cite{wei2022chain}. Nevertheless, CoT remains a cornerstone technique for enhancing the reasoning fidelity of LLMs and is particularly relevant to retrieval-based systems, where the ability to reason over retrieved evidence is paramount.
